\section{The continuum contradiction}
\label{sec:conclusion}

The path-pushing argument supplies the last property needed for a short
topological contradiction.  We first record why covering dimension one
allows a cohomology class on a closed subset to extend to the whole
space.

\begin{lemma}\label{end:restriction}
If $Z$ is a compact metric space with $\dim Z\le1$, $A\subset Z$ is
closed, and $k$ is an abelian coefficient group, then restriction induces
a surjection
\[
 \check H^1(Z;k)\longrightarrow\check H^1(A;k).
\]
\end{lemma}

\begin{proof}
Represent a class on $A$ by a cocycle on the nerve of a finite relative
open cover $\{V_i\}$ of $A$.  Choose open sets $W_i\subset Z$ with
$W_i\cap A=V_i$.  The cover consisting of the $W_i$ and $Z\setminus A$
has a finite open refinement $\mathcal U$ of multiplicity at most two.
Its nerve is a graph.  Retain the indices of the cover members when
restricting to $A$, even if two restrictions coincide. After empty
restrictions are discarded, the nerve of $\mathcal U|_A$ is a subgraph
of this graph. A refinement map pulls
the original cocycle back to that subgraph.  Extend the resulting
$1$-cochain by zero on all remaining edges of $N(\mathcal U)$.  Since
there are no $2$-simplices, the extension is a cocycle.  Its \v{C}ech
class restricts to the prescribed class on $A$.
\end{proof}

\begin{lemma}\label{end:paths-in-regions}
Let $L$ be a locally connected compact metric space with metric $d$, and let $U\subset L$
be connected and open. Any two points $a,b\in U$ are joined by a
continuous path in $U$.
\end{lemma}

\begin{proof}
Every connected open set $V\subset L$ is covered by connected open
sets $W$ whose closures lie in $V$ and whose diameters are at most any
prescribed $\varepsilon>0$. Indeed, choose a sufficiently small metric
ball about each point with closure in $V$, then a connected open
neighborhood inside a still smaller ball. The intersection graph of
this cover is connected: otherwise the unions belonging to its graph
components would separate $V$. Thus any two points of $V$ can be
joined by a finite chain $W_1,\ldots,W_m$ from this cover, with the
first point in $W_1$, the last in $W_m$, and
$W_i\cap W_{i+1}\ne\varnothing$. We may arrange $m\ge2$ by repeating
a member if necessary.

Starting with $[0,1]$, the set $U$, and endpoint values $a,b$, construct
successive finite interval partitions $\mathcal P_n$ as follows.
Associate to every interval $I\in\mathcal P_n$ a connected open set
$V_I$ containing its two already assigned endpoint values. For $n\ge1$
require $\operatorname{diam}V_I\le2^{-n}$. Inside $V_I$, choose the
finite chain above with closures in $V_I$ and diameters at most
$2^{-(n+1)}$. Subdivide $I$ into $m$ equal intervals and associate these
to the chain members. At each new division point assign a point of
the corresponding consecutive intersection. Keep the old endpoint
values. These assignments agree at shared endpoints, and both
endpoints of every child interval lie in its associated open set.

Every subdivision has at least two children, so
$\operatorname{mesh}(\mathcal P_n)\le2^{-n}$. The set $D$ of all
partition endpoints is therefore dense in $[0,1]$, and the compatible
assignments define $f:D\to L$. Nesting ensures that
$f(D\cap I)\subset V_I$ for every $I\in\mathcal P_n$.
For fixed $n\ge1$, let $\delta_n>0$ be the minimum length of a cell
of $\mathcal P_n$. If $s,t\in D$ and $|s-t|<\delta_n$, they lie in
the same or adjacent cells. In the adjacent case use the assigned
value at their common endpoint. In either case,
\[
 d(f(s),f(t))\le2^{1-n}.
\]
Thus $f$ is uniformly continuous and extends continuously to $[0,1]$
by completeness of $L$. For each cell $I\in\mathcal P_1$, its extended
image lies in $\overline{V_I}\subset U$. Hence the entire path lies in
$U$, and its endpoints are $a,b$.
\end{proof}

The image of a continuous path in a metric space contains an embedded
arc joining any two distinct points of that image; see
\cite[pp.~301--304]{Kamiya1930} for an elementary proof of Moore's
theorem. Consequently every connected open subset of $L$ in
Lemma~\ref{end:paths-in-regions} is arcwise connected.

\begin{lemma}\label{end:continuum}
Every nondegenerate locally connected compact metric continuum $L$
with
\[
 \dim L\le1,\qquad \check H^1(L;\mathbb F_2)=0
\]
has a cut point.
\end{lemma}

\begin{proof}
The continuum contains no embedded circle.  Such a circle $C$ would be
closed, and Lemma~\ref{end:restriction} would give a surjection from
$\check H^1(L;\mathbb F_2)=0$ onto
$\check H^1(C;\mathbb F_2)\cong\mathbb F_2$.

A locally connected compact metric continuum is a Peano continuum.
Lemma~\ref{end:paths-in-regions} and Moore's theorem give an arc
between any two distinct points of a connected open subset of $L$.
Choose an arc between two distinct points $a,b$ of
$L$ and an interior point $p$ of that arc.  If $p$ were not a cut point,
the connected open set $L\setminus\{p\}$ would contain a second arc
between $a$ and $b$.  The component of the first arc minus the second
that contains $p$ is an open subarc with both endpoints on the second
arc.  Its closure, together with the subarc of the second joining those
endpoints, is an embedded circle.  This contradiction shows that $p$
is a cut point.
\end{proof}

\begin{proof}[Proof of Theorem~\ref{thm:main}]
Suppose that $b_1^{(2)}(\widetilde Q)>0$.  The finite-energy construction,
the boundary topology proved in Propositions~\ref{topo:dimension}
and~\ref{topo:cech}, and the reduction of Proposition~\ref{block:setting}
give a nondegenerate compact metric continuum $L$ without cut points,
with $\dim L\le1$ and $\check H^1(L;\mathbb F_2)=0$.
Proposition~\ref{push:local-connected} makes $L$ locally connected.
Lemma~\ref{end:continuum} is a contradiction.  Hence
$b_1^{(2)}(\widetilde Q)=0$.

Lemma~\ref{found:reduction} gives $b_0^{(2)}=b_4^{(2)}=0$,
$b_3^{(2)}=b_1^{(2)}$, and vanishing in degrees above four.  Its
duality argument includes nontrivial orientation characters.
Consequently $b_p^{(2)}(\widetilde Q;\mathcal N(\Gamma))=0$ for every
$p\ge0$ with $p\ne2$, as asserted.
\end{proof}
